\section{Moldeado de recompensas y asignación de crédito temporal}

\subsection{LSTM convolucional y reward shaping}

\begin{figure}[H]
\centering
\includegraphics[width=0.85\textwidth]{slide-24.jpg}
\caption{Ejemplo de reward shaping que muestra la Lecture slide 24: se usa intermediate reward guidance en una multi-stage task para acortar el tiempo de entrenamiento.}
\end{figure}

El curso enfatiza que, en una multi-stage/hierarchical task, proporcionar un intermediate reward puede acelerar de manera significativa la convergencia del critic. La manera de hacerlo es descomponer el reward en \textit{completion reward} + \textit{auxiliary reward}, y usar dynamic weighting para controlar las gradient magnitudes.

\begin{knowledgebox}{Los guardrails de reward shaping}
El auxiliary reward debe mantenerse potential-based (es decir, no debe cambiar la optimal policy), de lo contrario incentivará al agent a desviarse del objective real. Por lo general se introduce una función de shaping $F(s,a,s')$ y se asegura que la suma $\sum_t \gamma^t F(s_t,a_t,s_{t+1})$ sea cero.
\end{knowledgebox}

\subsection{Planificador diferenciable y credit assignment}

El ponente mencionó un trick práctico: en control complejo se usa un differentiable planner para estimar la trajectory futura, y luego se usa el soft value del planner como señal de supervisión del critic. Esto puede reducir de forma significativa el drift del long-term credit assignment.

\begin{warningbox}{El soft planner requiere cómputo adicional}
Introducir differentiable planning aumenta el costo de inferencia en cada paso; es necesario asegurarse de que ese planner pueda ejecutarse en un entorno de low latency, de lo contrario arruinará el real-time control.
\end{warningbox}

\subsection{Resumen de la sección}

Reward shaping + la señal del planner es una solución realista para manejar multi-stage task y long-horizon credit assignment, pero hay que asegurarse de que el shaping se mantenga potential-based, sin acaparar el inference budget.
